% from ega1-10-fr.tex
\oldpage[I]{192}

sont commutatifs pour $m\leq n$~; autrement dit, $(f_n)$ est un
\emph{système inductif} de morphismes.
\end{env}

\begin{env}[10.6.8]
\label{I.10.6.8-fr}
Inversement, soit $(X_n)$ (resp. $(Y_n)$) un système inductif de
préschémas (usuels) satisfaisant aux conditions b) et c) de
(\hyperref[I.10.6.3-fr]{10.6.3}), et soit $\mathfrak{X}$ (resp.
$\mathfrak{Y}$) sa limite inductive. Par définition des limites
inductives, toute suite $(f_n)$ de morphismes $X_n\to Y_n$ formant un
système inductif admet une \emph{limite inductive}
$f:\mathfrak{X}\to\mathfrak{Y}$, qui est l'unique morphisme de
préschémas formels rendant commutatifs les diagrammes
\[
  \xymatrix{
    X_n\ar[r]^{f_n}\ar[d] &
    Y_n\ar[d]\\
    \mathfrak{X}\ar[r]_f &
    \mathfrak{Y}
  }
\]
\end{env}

\begin{proposition}[10.6.9]
\label{I.10.6.9-fr}
Soient $\mathfrak{X}$, $\mathfrak{Y}$ deux préschémas formels
localement noethériens, $\mathscr{J}$ (resp. $\mathscr{K}$) un Idéal
de définition de $\mathfrak{X}$ (resp. $\mathfrak{Y}$)~;
l'application $f\mapsto(f_n)$ définie dans
(\hyperref[I.10.6.7-fr]{10.6.7}) est une bijection de l'ensemble des
morphismes $f:\mathfrak{X}\to\mathfrak{Y}$ tels que
$f^*(\mathscr{K})\mathscr{O}_{\mathfrak{X}}\subset\mathscr{J}$, sur
l'ensemble des suites $(f_n)$ de morphismes rendant commutatifs les
diagrammes (\hyperref[I.10.6.7.1-fr]{10.6.7.1}).
\end{proposition}

\begin{proof}
Si $f$ est la limite inductive d'une telle suite, il faut montrer que
$f^*(\mathscr{K})\mathscr{O}_{\mathfrak{X}}\subset\mathscr{J}$. La
question étant locale sur $\mathfrak{X}$ et $\mathfrak{Y}$, on peut se
borner au cas où $\mathfrak{X}=\operatorname{Spf}(A)$,
$\mathfrak{Y}=\operatorname{Spf}(B)$ sont affines, $A$ et $B$ étant
adiques noethériens,
$\mathscr{J}=\mathfrak{J}^\Delta$,
$\mathscr{K}=\mathfrak{K}^\Delta$, où $\mathfrak{J}$ (resp.
$\mathfrak{K}$) est un idéal de définition de $A$ (resp. $B$). On a
alors $X_n=\operatorname{Spec}(A_n)$,
$Y_n=\operatorname{Spec}(B_n)$, avec
$A_n=A/\mathfrak{J}^{n+1}$ et $B_n=B/\mathfrak{K}^{n+1}$, en vertu de
(\hyperref[I.10.3.6-fr]{10.3.6}) et
(\hyperref[I.10.3.2-fr]{10.3.2})~;
$f_n=({}^a\varphi_n,\widetilde{\varphi_n})$, où les homomorphismes
$\varphi_n:B_n\to A_n$ forment un système projectif, donc
$f=({}^a\varphi,\widetilde{\varphi})$, où
$\varphi=\varprojlim\varphi_n$. La commutativité du diagramme
(\hyperref[I.10.6.7.1-fr]{10.6.7.1}) pour $m=0$ donne alors la
condition
$\varphi_n(\mathfrak{K}/\mathfrak{K}^{n+1})
\subset\mathfrak{J}/\mathfrak{J}^{n+1}$ pour tout $n$, donc, en
passant à la limite projective, $\varphi(\mathfrak{K})\subset
\mathfrak{J}$, et cela entraîne
$f^*(\mathscr{K})\mathscr{O}_{\mathfrak{X}}\subset\mathscr{J}$
(\hyperref[I.10.5.6-fr]{10.5.6, (ii)}).
\end{proof}

\begin{corollary}[10.6.10]
\label{I.10.6.10-fr}
Soient $\mathfrak{X}$, $\mathfrak{Y}$ deux préschémas formels
localement noethériens, $\mathscr{T}$ le plus grand Idéal de définition
de $\mathfrak{X}$ (\hyperref[I.10.5.4-fr]{10.5.4}).
\begin{enumerate}
  \item[\rm{(i)}] Pour tout Idéal de définition $\mathscr{K}$ de
  $\mathfrak{Y}$ et tout morphisme
  $f:\mathfrak{X}\to\mathfrak{Y}$, on a
  $f^*(\mathscr{K})\mathscr{O}_{\mathfrak{X}}\subset\mathscr{T}$.
  \item[\rm{(ii)}] Il y a correspondance biunivoque canonique entre
  $\operatorname{Hom}(\mathfrak{X},\mathfrak{Y})$ et l'ensemble des
  suites $(f_n)$ de morphismes rendant commutatifs les diagrammes
  (\hyperref[I.10.6.7.1-fr]{10.6.7.1}), où
  $X_n=(\mathfrak{X},
  \mathscr{O}_{\mathfrak{X}}/\mathscr{T}^{n+1})$,
  $Y_n=(\mathfrak{Y},
  \mathscr{O}_{\mathfrak{Y}}/\mathscr{K}^{n+1})$.
\end{enumerate}
\end{corollary}

\begin{proof}
(ii) résulte aussitôt de (i) et de
(\hyperref[I.10.6.9-fr]{10.6.9}). Pour démontrer (i), on peut se borner
au cas où $\mathfrak{X}=\operatorname{Spf}(A)$,
$\mathfrak{Y}=\operatorname{Spf}(B)$, $A$ et $B$ étant noethériens,
$\mathscr{T}=\mathfrak{T}^\Delta$,
$\mathscr{K}=\mathfrak{K}^\Delta$, où $\mathfrak{T}$ est le plus grand
idéal de définition de $A$ et $\mathfrak{K}$ un idéal de définition de
$B$. Soit $f=({}^a\varphi,\widetilde{\varphi})$, où
$\varphi:B\to A$ est un homomorphisme continu~; comme les éléments de
$\mathfrak{K}$ sont topologiquement nilpotents
(\hyperref[0.7.1.4-fr]{0, 7.1.4, (ii)}), il en est de même de ceux de
$\varphi(\mathfrak{K})$, donc
$\varphi(\mathfrak{K})\subset\mathfrak{T}$ puisque $\mathfrak{T}$ est
l'ensemble des éléments topologiquement nilpotents de $A$
(\hyperref[0.7.1.6-fr]{0, 7.1.6})~; d'où la conclusion en vertu de
(\hyperref[I.10.5.6-fr]{10.5.6, (ii)}).
\end{proof}

\begin{corollary}[10.6.11]
\label{I.10.6.11-fr}
Soient $\mathfrak{S}$, $\mathfrak{X}$, $\mathfrak{Y}$ trois
préschémas formels localement noethériens,
$f:\mathfrak{X}\to\mathfrak{S}$,
$g:\mathfrak{Y}\to\mathfrak{S}$ des morphismes faisant de
$\mathfrak{X}$ et $\mathfrak{Y}$ des $\mathfrak{S}$-préschémas
formels. Soit $\mathscr{J}$ (resp. $\mathscr{K}$, $\mathscr{L}$) un
Idéal de définition de $\mathfrak{S}$ (resp. $\mathfrak{X}$,
$\mathfrak{Y}$), et supposons que
$f^*(\mathscr{J})\mathscr{O}_{\mathfrak{X}}\subset\mathscr{K}$,
$g^*(\mathscr{J})\mathscr{O}_{\mathfrak{Y}}=\mathscr{L}$~; posons
$S_n=(\mathfrak{S},
\mathscr{O}_{\mathfrak{S}}/\mathscr{J}^{n+1})$,
$X_n=(\mathfrak{X},
\mathscr{O}_{\mathfrak{X}}/\mathscr{K}^{n+1})$,
$Y_n=(\mathfrak{Y},
\mathscr{O}_{\mathfrak{Y}}/\mathscr{L}^{n+1})$. Il y a alors
correspondance

